\section{Divisor arithmetic without factoring supports}
\label{sec:divisors}

The lifting argument uses rational divisors, although the finite points in
their supports are not given individually.  We now give a representation in
which addition, reduction, and evaluation require only polynomial arithmetic
and linear algebra.  Two further operations will be needed: summing divisors
computed simultaneously at the unknown roots of $F$, and evaluating a
principal function whose divisor has very large coefficients at infinity.
Linear-algebra algorithms for Riemann--Roch spaces and divisor arithmetic
are established tools; see Hess~\cite{Hess02} and
Khuri-Makdisi~\cite{KhuriMakdisi04}.  We give the representation and size
bounds explicitly because both the genus and the degree of the cover vary
in this application.

Throughout this section, $k$ is a finite field, $q$ is an odd prime different
from its characteristic, $k$ contains a specified primitive $q$th root
$\zeta$, and $F\in k[X]$ is monic and square-free of degree $N\geq q$, with
$q\mid N$.  We use the curve, automorphism, and rational points at infinity
introduced above.  Thus
\[
 C:\ Y^q=F(X),\qquad
 \mathcal L=k(C),\qquad
 O=k[X,Y]/(Y^q-F(X)),\qquad
 g=\frac{(q-1)(N-2)}2.
\]
The automorphism $\sigma$ sends points by $Y\mapsto\zeta Y$ and acts on
functions by $Y\mapsto\zeta^{-1}Y$.  All divisors in the computation are
$k$-rational divisors, written as sums of closed points.  For a divisor $D$,
write
\[
 \|D\|=\sum_P |v_P(D)|[k(P):k]
\]
for its absolute degree.  Bounds polynomial in $\|D\|$ do not mean bounds
polynomial in the bit length of arbitrary divisor coefficients.  Large
coefficients supported at infinity are treated separately at the end of
this section.

\begin{proposition}[Effective divisor arithmetic]\label{div:arithmetic}
On the curve $C$ above, the following operations are deterministic and use
a number of field operations polynomial in $q,N$, the degrees occurring in
the input representation, the absolute degrees of the divisors
materialized during the operation, and the parameter $\ell$ in the third
operation:
\begin{enumerate}
 \item divisor addition, negation, coefficientwise minimum and maximum,
 and the action of $\sigma$;
 \item construction of the divisor of an explicitly represented nonzero
 function, and evaluation of its valuation and leading coefficient at a
 specified ramification point or point at infinity;
 \item computation of the Riemann--Roch space
 $L(D+\ell\infty_0)$ for $\ell\geq 0$;
 \item replacement of a degree-zero divisor $D$ by a linearly equivalent
 divisor $D'$ satisfying $\|D'\|\leq 2g$, with a function $a$ such that
 $D'=D+\operatorname{div}(a)$.
\end{enumerate}
These algorithms do not require factorization of the polynomials defining
the finite supports.  Their decisions use only field zero tests, so they
admit simultaneous execution as in Lemma~\ref{ff:simultaneous}.
\end{proposition}

We prove the proposition by constructing the algorithms.  The final two
subsections extend them to simultaneous sums and to binary infinity
coefficients.

\subsection{Fractional ideals and finite linear algebra}

Smoothness of the affine curve implies that $O$ is a Dedekind domain.  It is
a free $k[X]$-module with basis $1,Y,\ldots,Y^{q-1}$.  The finite part of a
divisor $D$ is represented by the fractional ideal
\begin{equation}\label{div:ideal}
 I(D)=\{a\in\mathcal L:v_P(a)\geq v_P(D)
                 \text{ at every finite point }P\}.
\end{equation}
The zero function is included in this definition.  The rest of the
representation is the integer vector
$(v_{\infty_0}(D),\ldots,v_{\infty_{q-1}}(D))$.

To store a fractional ideal $I$, choose a monic polynomial $H\in k[X]$
such that
\begin{equation}\label{div:bound}
 HO\subseteq I\subseteq H^{-1}O.
\end{equation}
Store the image of $I$ as a $k$-subspace of $H^{-1}O/HO$.  Multiplication
by $H$ identifies this quotient with $O/H^2O$, of dimension $2q\deg H$.
Thus a subspace basis consists of polynomial vectors modulo $H^2$ in the
standard $Y$-basis.  Lifting this basis and adjoining $H$ gives a finite
list of $O$-generators of $I$.  Conversely, from such generators and a
bound~\eqref{div:bound}, take their images and repeatedly enlarge their
$k$-span under multiplication by $X$ and $Y$.  The process stabilizes in
at most $2q\deg H$ strict enlargements and gives the required ideal image.
The convention $H=1$ represents $I=O$ with a zero-dimensional quotient.

The identities
\begin{align*}
 I(D+E)&=I(D)I(E),& I(-D)&=I(D)^{-1},\\
 I(\min(D,E))&=I(D)+I(E),&
 I(\max(D,E))&=I(D)\cap I(E)
\end{align*}
reduce the first group of operations to finite linear algebra.  For a
product, take pairwise products of generators and use the bound $H_DH_E$.
For a sum or intersection, embed both subspaces using that same common
bound, and take their sum or intersection.  For inversion, retain the
bound $H$: an element $a\in H^{-1}O$ belongs to $I^{-1}$ precisely when
$au\in O$ for every generator $u$ of $I$.  These are linear conditions
on $H^{-1}O/HO$.  They are independent of the chosen representative of $a$,
since $HI\subseteq O$.  The action of $\sigma$ is substitution of
$\zeta^{-1}Y$ for $Y$ in the generators.  The corresponding operations on
the infinity vectors are immediate.

A known ramification point $P=(x,0)$ has ideal $(X-x,Y)$, with bound
$H=X-x$.  For a nonzero function $a$, arithmetic in $\mathcal L$ uses
rational coefficients in $k(X)$ in the standard $Y$-basis.  Inversion is
linear algebra over $k(X)$.  Clearing denominators in both $a$ and $a^{-1}$
gives a polynomial $H$ for which $aO$ satisfies~\eqref{div:bound}; hence
the finite part of $\operatorname{div}(a)$ is computable without finding
any of its points.

Intermediate bound polynomials need not be minimal.  The following
compression both controls their degrees and supplies the norm data used
in class division.

\begin{lemma}[Compression and pushforward]\label{div:compression}
From the representation of $I=I(D)$ one can compute monic polynomials
$h^+,h^-\in k[X]$ satisfying
\[
 \deg h^+=\sum_{P\text{ finite}}\max(v_P(D),0)[k(P):k],\qquad
 \deg h^-=\sum_{P\text{ finite}}\max(-v_P(D),0)[k(P):k].
\]
The polynomial $H=h^+h^-$ satisfies~\eqref{div:bound}.  If $\deg D=0$ and
$j=h^+/h^-$, then
\begin{equation}\label{div:orbitnorm}
 \operatorname{div}_C(j)=\sum_{a=0}^{q-1}\sigma^aD.
\end{equation}
\end{lemma}

\begin{proof}
Form the integral ideals
$I^+=I\cap O$ and $I^-=I^{-1}\cap O$.  Let $h^+$ and $h^-$ be the
characteristic polynomials of multiplication by $X$ on $O/I^+$ and
$O/I^-$, respectively.  These quotients are computed from the stored
subspaces; the original bound $H_0$ has $H_0O\subseteq I^+,I^-$, so it
suffices to work in $O/H_0O$.  A quotient at a closed point $P$ of
multiplicity $d$ has a filtration with $d$ successive residue fields
$k(P)$.  Its $k$-dimension is therefore $d[k(P):k]$, proving the degree
formulas.

Cayley--Hamilton gives $h^+O\subseteq I^+\subseteq I$ and
$h^-O\subseteq I^-\subseteq I^{-1}$.  The second inclusion implies
$I\subseteq(h^-)^{-1}O$.  Both sides of~\eqref{div:bound} now follow
for $H=h^+h^-$.  Recompute the ideal subspace with this bound after each
operation.  In particular, its dimension is at most twice $q$ times the
absolute degree of the finite divisor.

For the last assertion, let $Q=X(P)$ be a finite point downstairs and
$f_Q$ its monic defining polynomial.  On $k(P)$, multiplication by $X$
has characteristic polynomial $f_Q^{[k(P):k(Q)]}$.  Applying this fact to
the preceding filtrations shows that the finite divisor of $j$ on
$\mathbb P^1$ is $X_*D$ at the finite points.  Since $\deg D=0$, the
infinity coefficient agrees as well.  Finally, for this cyclic cover,
$X^*X_*D=\sum_{a=0}^{q-1}\sigma^aD$, including the ramified points and
points with nontrivial residue extension.  This proves
\eqref{div:orbitnorm}.
\end{proof}

\subsection{Local expansions and Riemann--Roch reduction}

The ideal representation controls all finite points collectively.  At the
few specified points where a value or inequality is required, local
expansions give the additional information.

At infinity use $z=X^{-1}$ and $U=Yz^{N/q}$.  The equation
\[
 U^q=z^NF(z^{-1}),\qquad U(0)=\zeta^j
\]
determines the expansion at $\infty_j$ by simple-root Hensel lifting.
At $P=(x,0)$ use $Y$ as uniformizer and solve $F(X)=Y^q$ with $X(0)=x$;
the derivative $F'(x)\ne0$ again permits coefficient-by-coefficient
lifting.  The valuation and leading coefficient of a nonzero function
are then found by substitution and division of the leading terms.

Here is an explicit precision bound.  Write a function as
$A(X,Y)/B(X)$, where $\deg_Y A<q$ and every coefficient in $X$, as well
as $B$, has degree at most $E$.  The polynomial $A$ has poles only at
infinity, with total pole degree at most $q(E+N)$; the same bound holds
for $B$.  The degree of its zero divisor equals its pole degree, so no
nonzero numerator or denominator can have an order of vanishing greater
than this bound at a rational point.  Consequently local substitutions
through precision
\begin{equation}\label{div:precision}
 8q(E+N+1)
\end{equation}
suffice to identify their first nonzero Laurent coefficients, even after
the possible pole at infinity has been removed.  For a represented
divisor, its coefficient at a specified finite point is the minimum of
the valuations of a list of ideal generators; zero generators may be
discarded.

We next turn Riemann--Roch existence into a linear computation.  Recall
that
\[
 L(V)=\{a\in\mathcal L:\operatorname{div}(a)+V\geq0\}\cup\{0\}
\]
is a finite-dimensional $k$-vector space.  Given $D$ and $\ell\geq0$,
choose a bound $H$ for $I(-D)$.  Every function in $L(D+\ell\infty_0)$
has the form
\begin{equation}\label{div:rransatz}
 a=H^{-1}\sum_{j=0}^{q-1}a_j(X)Y^j,
 \qquad \deg a_j\leq \deg H+\|D\|+\ell.
\end{equation}
To justify the degree bound, put
\[
 M=\max_{a_j\ne0}\bigl(\deg a_j+jN/q-\deg H\bigr).
\]
The coefficient of $z^{-M}$ in the expansions at infinity is a nonzero
polynomial of degree less than $q$ evaluated at the $q$ distinct values
$1,\zeta,\ldots,\zeta^{q-1}$.  It cannot vanish at all of them.  Thus
at one infinity point the pole order, interpreted as $-v_{\infty_j}(a)$,
is exactly $M$.  The defining inequalities for $L(D+\ell\infty_0)$ give
$M\leq\|D\|+\ell$, which implies~\eqref{div:rransatz}.

In this finite ansatz, membership in $I(-D)$ is a subspace condition, and
\[
 v_{\infty_j}(a)\geq-v_{\infty_j}(D)-\ell\,\mathbf1_{j=0}
 \qquad(0\leq j<q)
\]
is a list of linear conditions on Laurent coefficients.  Expansions to
the bound~\eqref{div:precision}, increased by
$8q(\|D\|+\ell+1)$, suffice.  Solving the resulting linear system
therefore computes the whole space $L(D+\ell\infty_0)$.

For degree-zero $D$, the Riemann--Roch inequality
\cite[Tag~0BS6]{StacksProject} gives
$\dim_kL(D+g\infty_0)\geq1$.  Choose a nonzero solution $a$ by a fixed
pivot convention and define
\begin{equation}\label{div:reduce}
 \operatorname{ReduceDivisor}(D)=D'=D+\operatorname{div}(a).
\end{equation}
Since $D'+g\infty_0$ is effective of degree $g$, we have
$\|D'\|\leq2g$.  Retain $a$ whenever a later computation needs the
linear equivalence.  If a degree-zero divisor $D$ is known to be
principal, taking $\ell=0$ instead produces a nonzero $a$ with
$D+\operatorname{div}(a)=0$: the effective divisor on the left has degree
zero.

All vector spaces and polynomial degrees in these constructions are
polynomial in the parameters in Proposition~\ref{div:arithmetic}.
Dense polynomial operations, fraction reduction over $k(X)$, and linear
algebra therefore prove its complexity assertion.  For example, the
ideal quotient has dimension $2q\deg H$, and the ansatz
\eqref{div:rransatz} has at most
$q(\deg H+\|D\|+\ell+1)$ unknowns.  Rational linear algebra does not
cause uncontrolled degree growth: after common denominators are cleared,
minors of an $s\times s$ matrix with polynomial entries of degree at
most $d$ have degree at most $sd$.  Reduced elimination entries and
solutions are ratios of such minors.  This completes the proof of
Proposition~\ref{div:arithmetic}.

\subsection{Summing divisors given over a product algebra}

The odd-degree splitter first produces one divisor for each root of $F$,
through a computation over $A=k[T]/F(T)$.  The next lemma computes their
sum over $k$, while the roots remain unknown.  A norm from the finite
algebra $A$ will replace the unavailable product over its components.

\begin{lemma}[Simultaneous divisor sum]\label{div:trace}
Suppose $A\simeq k^s$ is an explicitly represented split algebra.  Let
$I_1,\ldots,I_s$ be fractional ideals of $O$ supplied together over $A$,
with monic bounds $H_i$ and a common number $b$ of generators
$u_{i0},\ldots,u_{i,b-1}$, including $H_i$.  If the prime subfield
contains more than $s(b-1)$ elements, their product can be computed by
polynomial arithmetic, determinants, and ideal closure, without a
decomposition of $A$.  Together with addition of the infinity vectors,
this computes the sum of the corresponding divisors.
\end{lemma}

\begin{proof}
Let $H\in A[X]$ denote the tuple of bounds and let $u_j$ denote the tuple
of $j$th generators.  Form
\[
 H_* =\operatorname{Norm}_{A/k}(H)=\prod_{i=1}^s H_i,
 \qquad
 w_a=\operatorname{Norm}_{A/k}\left(\sum_{j=0}^{b-1}a^ju_j\right)
     =\prod_{i=1}^s\left(\sum_{j=0}^{b-1}a^ju_{ij}\right).
\]
The norm is the determinant of multiplication after extending scalars
to the relevant polynomial or function ring.  The displayed products
are identities proving correctness, not instructions to recover the
components.  Denominators are cleared by multiplying the $i$th component
by $H_i$; the resulting determinant lies in $O$, and division by $H_*$
recovers $w_a$.  Polynomial determinants can be computed before reducing
by $Y^q-F(X)$.  Use $s(b-1)+1$ distinct values of $a$ from the prime
subfield.

Each $w_a$ belongs to $\prod_i I_i$.  Fix a finite closed point $P$ and
write $d_i=\min_jv_P(u_{ij})$.  After dividing the $i$th combination by
a local uniformizer to power $d_i$, its residue is a nonzero polynomial
in $a$, over $k(P)$, of degree at most $b-1$.  At most $b-1$ of our
values fail to attain valuation $d_i$.  Thus some chosen $a$ attains all
$d_i$ simultaneously, and $v_P(w_a)=\sum_i d_i$.  Fractional ideals in
a Dedekind domain are determined by these minimum valuations.  The
$w_a$ therefore generate exactly $\prod_i I_i$.  The polynomial $H_*$
is a valid bound, so finite ideal closure computes its stored
representation.

Infinity coefficients are integers.  In a simultaneous field execution
they are the same in every surviving component: any different outcome
of a scalar zero test would already have supplied a factor.  Their sum
is therefore obtained by multiplying the common vector by $s$.  More
generally, if the computation has retained separate pieces, sum their
vectors with the corresponding dimensions.
\end{proof}

For the application $s=N$.  The number $b$ can be bounded by
$2q\deg H_i+1$ in the compressed representation, and the characteristic
cutoff in Section~\ref{sec:complexity} guarantees the required distinct
scalars.  Matrix determinants and ideal closure have sizes polynomial
in these parameters.

\subsection{Binary coefficients at infinity and function circuits}

Riemann--Roch reduction has cost polynomial in absolute degree.  It must
not be applied directly to an infinity divisor whose coefficients are
exponentially large integers.  We instead perform the class arithmetic
by binary additions and retain the accompanying rational function as a
circuit.

\begin{lemma}[Principal functions with large infinity coefficients]
\label{div:circuits}
Let $D$ be a represented degree-zero divisor of bounded absolute degree,
and let $W=\sum_jw_j\infty_j$ have degree zero, with integer coefficients
given in binary.  If
$Z=(1-\sigma)D-W$ is principal, a circuit for a function $h$ satisfying
$\operatorname{div}(h)=Z$ can be computed using a number of operations
polynomial in $q,N,\|D\|$, the representation size, and
$\max_j\log(2+|w_j|)$.  The circuit uses multiplication and inversion
of explicitly represented small functions.  At each known ramification
point $P$, it computes the nonzero value $h(P)$ within the same type of
bound, without expanding $h$.
\end{lemma}

\begin{proof}
Write $W=\sum_{j=1}^{q-1}w_j(\infty_j-\infty_0)$.  Build the divisor
$Z$ from $(1-\sigma)D$ and these signed multiples by binary doubling
and addition.  At every stage retain a reduced divisor $D_0$ and a
function circuit $w$ with the invariant
\begin{equation}\label{div:circuitinvariant}
 D_0=D_{\mathrm{target}}+\operatorname{div}(w).
\end{equation}
On adding two stages, reduce $D_1+D_2$ with reducing function $u$ and
replace the function circuit by $w_1w_2u$.  For doubling it becomes
$w_1^2u$; negation uses inversion followed, if needed, by reduction.
After the initial reduction, each materialized sum has absolute degree
at most $4g$, and the number of reductions is
$O(q\max_j\log(2+|w_j|))$, in addition to the initial operations on $D$.

At the end obtain a small divisor $\widetilde Z=Z+\operatorname{div}(w)$.
It is principal, so the $\ell=0$ Riemann--Roch computation gives $a$ with
$\widetilde Z+\operatorname{div}(a)=0$.  Consequently
\[
 h=(aw)^{-1}
\]
has divisor $Z$.  All leaves of this circuit are functions produced by
reductions on small divisors and hence have polynomial representation
size.

At $P$, expand each leaf and store its integer valuation and its nonzero
leading coefficient.  Multiplication adds valuations and multiplies
leading coefficients; inversion negates valuations and inverts leading
coefficients.  The resulting integers have bit length polynomial in the
circuit length and leaf bounds.  Because $P$ is fixed by $\sigma$ and
$W$ is supported at infinity, $v_P(Z)=0$.  The final valuation is
therefore zero, and the final leading coefficient is exactly $h(P)$.
This also explains why the intermediate objects are leading coefficients
of Laurent series, rather than values of functions that may have poles.
\end{proof}
